% =====================================================================
%  c21 —— 原文印刷页 166–171（PDF 第 167–172 页）
%  附录 A 的开头：A. Symplectic Geometry（A.1 Kähler quotients 的
%  A.1.1 The moment map、A.1.2 Symplectic Quotient、A.1.3 Kähler quotient）
%  与 B. Lightning Review Of Kähler Quantization 的 B.1 The Bargmann
%  Representation。本片段无插图。
%  main.tex 已在本片段之前发出 \appendix，故此处只用 \section/\subsection/
%  \subsubsection，编号 A、A.1、A.1.1 由 LaTeX 自动给出。
% =====================================================================

\section{辛几何} % 原文: A. Symplectic Geometry

\begin{keytakeaways}
  \item 物理图像：相空间就是一个辛流形 $(M,\omega)$，即带有非退化二形式的流形；$\omega$ 的逆给出 Poisson 括号 (A.2)，于是经典力学的运动学完全编码在 $\omega$ 里，度规并不出现。
  \item 数学工具：Darboux 定理说明辛形式没有非平凡的局部不变量，辛结构的信息全是整体的；连续对称性则由矩映射 $\mu:M\to\mathfrak{g}^*$ 承载，(A.10) 把「每条生成元对应一个守恒荷」写成配对，(A.11) 要求它关于余伴随作用等变。
  \item 核心结论：辛商（Hamilton 约化）$M/\!/_\zeta G=\mu^{-1}(\zeta)/G$ (A.30) 给出约化后的相空间；而 Poisson 括号给出的 Lie 代数 2-上循环 $c^{AB}$ (A.26) 若代表非平凡上同调类，则同时施加第一类约束会遇到反常。
  \item 关键例子：$\mathbb{C}^N$ 上按权 $Q^a$ 的 $U(1)$ 作用，商空间的性质由 $Q^a$ 的可公度性、互素性与符号决定——权两两互素时商空间光滑，出现公共因子时留下 $\mathbb{Z}_r$ 轨形奇点，而权异号时约束面非紧；当 $Q_a=\pm1$ 时辛商是 $\mathcal{O}(1)$ 丛的全空间。
  \item 与场论的联系：超对称场论势能中的 $D$ 项恰好是靶空间上 $G$-作用的矩映射，矩映射的能级 $\zeta^A$ 就是 Fayet–Iliopoulos 项。
\end{keytakeaways}

回忆：相空间是一个辛流形（symplectic manifold），即带有非退化二形式
\begin{equation}
  \omega=\frac{1}{2}\omega_{\mu\nu}dx^{\mu}\wedge dx^{\nu}
  \tag{A.1}
  \label{eq:A.1}
\end{equation}
的流形；由此可以定义相空间上两个函数的 Poisson 括号：
\begin{equation}
  \{f_1,f_2\}=\omega^{\mu\nu}\frac{\partial f_1}{\partial x^{\mu}}\frac{\partial f_2}{\partial x^{\nu}}
  \tag{A.2}
  \label{eq:A.2}
\end{equation}

辛几何的第一个主要定理是：辛形式没有非平凡的局部不变量。这就是 Darboux 定理：它说在 $\mathcal{P}$ 上总存在局部坐标 $(p_i,q^i)$，使得
\begin{equation}
  \omega=\sum_i dp_i\wedge dq^i
  \tag{A.3}
  \label{eq:A.3}
\end{equation}
{\small\itshape ［原文标记：$\clubsuit$这并不难证明。要不要给出证明？$\clubsuit$］}
% [原文待办标记，原文 p.166 页边竖排；作者在此自问是否补上 Darboux 定理的证明，最终未补]

相空间是讨论经典力学的恰当舞台，也是通向量子化的良好出发点。

\subsection{Kähler 商：快速回顾} % 原文: A.1 Kahler quotients: A lightning review

\subsubsection{矩映射} % 原文: A.1.1 The moment map

设 $(M,\omega)$ 是一个相空间，即一个辛流形。

回忆：给定 $M$ 上的函数 $h$，可以由 $\iota(V)\omega=dh$ 作出与之对应的 Hamilton 向量场；用局部坐标写出就是
\begin{equation}
  V^i=\omega^{ij}\partial_j h
  \tag{A.4}
  \label{eq:A.4}
\end{equation}
% [Tutor Note: 符号约定易错。(A.2) 把 Poisson 括号定义为 $\{f,g\}=\omega^{ij}\partial_if\,\partial_jg$，而 (A.4) 取 $V^i=\omega^{ij}\partial_jh$，即 $\iota(V)\omega=dh$。直接代入会得到 $V_\mu(f)=-\{\mu,f\}=\{f,\mu\}$，与 (A.8) 第一行到第二行所默认的 $\{\mu,f\}=\xi^i\partial_if$ 差一个负号。核对时请把 (A.2)、(A.4)、(A.8) 三式并排验证（用 $\omega^{12}=1$、$\mu=x^1$、$f=x^2$ 试算即可看出）；作者本人在 p.171 的 (B.1) 旁也标了「SIGN???」，可见他对符号是随手处理的]

向量场 $V$ 生成一个保持辛形式的微分同胚群。

反过来，设 $(M,\omega)$ 是辛流形，且 Lie 群 $G$ 辛作用（acts symplectically）：
\begin{equation}
  g^{*}(\omega)=\omega\qquad\forall g\in G
  \tag{A.5}
  \label{eq:A.5}
\end{equation}

把这一点用到由 $T\in\mathfrak{g}$ 生成的单参数子群上，便得到向量场 $\xi=\xi(T)$，使得
\begin{align}
  0 &= \mathcal{L}_{\xi}\omega \notag\\
    &= \big(d\iota(\xi)+\iota(\xi)d\big)\omega \label{eq:A.6}\tag{A.6}\\
    &= d\big(\iota(\xi)\omega\big) \notag
\end{align}

因此，若 $H^{1}_{DR}(M)=0$，我们就可以断定
\begin{equation}
  \iota(\xi)\omega=d\mu
  \tag{A.7}
  \label{eq:A.7}
\end{equation}
其中 $\mu:M\to\mathbb{R}$ 是某个整体良定义的函数。注意 $\mu$ 只定义到相差一个常数。它依赖于 $T$，所以我们把它在一点处的值记作 $\mu(p;T)$，$p\in M$、$T\in\mathfrak{g}$。

在 Hamilton 动力学中，$\mu$ 就是与向量场 $\xi$ 所生成对称性相应的荷（charge）。它生成与该对称性相联系的变换：
\begin{align}
  \{\mu,f\} &= \omega^{ij}\partial_i\mu\,\partial_jf \notag\\
           &= \xi^i\partial_if \label{eq:A.8}\tag{A.8}\\
           &= \mathcal{L}_{\xi}(f) \notag
\end{align}
其中最后一行用到了关于 $\xi$ 的 Lie 导数。

对每个元素 $T\in\mathfrak{g}$，我们都有一个守恒荷 $\mu(\cdot;T)$。由于 $T\mapsto\xi(T)$ 关于 $T$ 线性，$T\mapsto\mu(\cdot;T)$ 也关于 $T$ 线性。以「关于 $T$ 线性」的方式对每个 $T\in\mathfrak{g}$ 给出一个函数 $M\to\mathbb{R}$，等价于给出单个函数
\begin{equation}
  \mu:M\to\mathfrak{g}^{*}
  \tag{A.9}
  \label{eq:A.9}
\end{equation}

这个函数称为矩映射（moment map）。按定义它对一切 $T\in\mathfrak{g}$ 满足
\begin{equation}
  \langle T,\mu(p)\rangle:=\mu(p;T)
  \tag{A.10}
  \label{eq:A.10}
\end{equation}
此外还有
\begin{equation}
  \mu(g\cdot p)=\operatorname{Ad}^{*}(g)\big(\mu(p)\big)
  \tag{A.11}
  \label{eq:A.11}
\end{equation}
其中右端是余伴随作用（coadjoint action）。
% [Tutor Note: 术语取舍。"moment map" 的中文译法有「矩映射」「动量映射」「矩量映射」三种；此处统一取「矩映射」，因为它保留 "moment"（而非 momentum）的字面，也与本文后面 $D$ 项、Fayet–Iliopoulos 项的语境相容。作者紧接着自己解释 "the word 'moment' is short for 'momentum'"，故正文首次出现处括注英文原词，读者可自行对照]

取 $\mathfrak{g}$ 的一组基 $T^A$，我们可以写
\begin{equation}
  \mu=\sum_A\mu^A T^{*}_A
  \tag{A.12}
  \label{eq:A.12}
\end{equation}
其中 $\mu^A=\mu(T^A)$ 是 $M$ 上的普通函数。

“moment” 一词是 “momentum” 的缩写。它可以是线动量，也可以是角动量，如下面的例子所示：

\textbf{例 1}：设 $M=T^{*}\mathbb{R}^n$，取 Darboux 坐标 $(\vec{q},\vec{p})=(\vec{q}^{\,i},\vec{p}_i)$，于是
\begin{equation}
  \omega=\sum_i dp_i\wedge dq^i
  \tag{A.13}
  \label{eq:A.13}
\end{equation}

设 $G=\mathbb{R}^n$ 通过平移作用：
\begin{equation}
  \vec{a}\cdot(\vec{q},\vec{p})=(\vec{q}+\vec{a},\vec{p})
  \tag{A.14}
  \label{eq:A.14}
\end{equation}

我们还有：作为向量空间 $\mathfrak{g}\cong\mathbb{R}^n$，其 Lie 括号恒为零，而无穷小作用由 $g(t)=e^{t\vec{v}}$ 给出：
\begin{equation}
  g(t):(\vec{q},\vec{p})\to(\vec{q}+t\vec{v},\vec{p})
  \tag{A.15}
  \label{eq:A.15}
\end{equation}

于是
\begin{equation}
  \xi(\vec{v})=v^i\frac{\partial}{\partial q^i}
  \tag{A.16}
  \label{eq:A.16}
\end{equation}

并且
\begin{equation}
  \iota\big(\xi(\vec{v})\big)\omega=-v^idp_i=d\mu(\vec{v})
  \tag{A.17}
  \label{eq:A.17}
\end{equation}
其中
\begin{equation}
  \mu(q^i,p_i;\vec{v})=-v^ip_i
  \tag{A.18}
  \label{eq:A.18}
\end{equation}

\textbf{例 2} 转向我们的主要例子：设 $G$ 线性作用在向量空间 $V=\{(x^1,\dots,x^N)\}$ 上，辛形式为 $\omega_{ij}dx^idx^j$，其中 $\omega_{ij}$ 是常值反对称矩阵。设 $\rho(T^A)^i{}_j$ 是表示该 Lie 代数的 $N\times N$ 矩阵。则
\begin{equation}
  \xi^A:=\xi(T^A)=x^j\rho(T^A)^i{}_j\frac{\partial}{\partial x^i}
  \tag{A.19}
  \label{eq:A.19}
\end{equation}

于是矩映射为
\begin{equation}
  \mu^A(x)=\frac{1}{2}x^j\rho(T^A)^i{}_j\omega_{ik}x^k
  \tag{A.20}
  \label{eq:A.20}
\end{equation}
% [Tutor Note: (A.20) 的因子 $1/2$ 来自「对 $x^j\omega_{ik}x^k$ 求微分后除以 2」：把 (A.20) 代入 $d\mu^A=\iota(\xi^A)\omega$，用到 $\omega_{ik}$ 反对称与紧接 (A.20) 之后那句「$\rho(T^A)^i{}_j\omega_{ik}$ 关于 $jk$ 对称」，两份贡献相加恰好抵消 $1/2$。若漏掉那条对称性条件，就会算出多一个因子 2]

（该群辛作用，当且仅当矩阵 $\rho(T^A)^i{}_j\omega_{ik}$ 关于 $jk$ 对称。）

特殊化到 $M=T^{*}\mathbb{R}^n$ 并取标准 Darboux 坐标的情形，我们可以考虑嵌入 $GL(n,\mathbb{R})\hookrightarrow Sp(2n,\mathbb{R})$，其中
\begin{equation}
  \begin{pmatrix}
    A & \\
      & A^{tr,-1}
  \end{pmatrix}
  \tag{A.21}
  \label{eq:A.21}
\end{equation}
% [Tutor Note: (A.21) 的记法 $A^{tr,-1}$ 应读作 $(A^{\mathrm{tr}})^{-1}$（等价于 $(A^{-1})^{\mathrm{tr}}$），不是「$A^{\mathrm{tr}}$ 减 1」；整个块矩阵是 $\mathrm{diag}\big(A,(A^{\mathrm{tr}})^{-1}\big)$，它保持 (A.13) 的 $\sum_i dp_i\wedge dq^i$。原文把转置与求逆压进一个上标，转写时若不还原，读者无法判断第二块到底是 $A^{\mathrm{tr}}$ 还是 $A^{-1}$]

对群元素 $g(t)=1+te_{ij}\in GL(n,\mathbb{R})$，我们有
\begin{equation}
  g(t):(\vec{q},\vec{p})\to(\vec{q}+tq_j\vec{e}_i,\ \vec{p}-tp_i\vec{e}_j)
  \tag{A.22}
  \label{eq:A.22}
\end{equation}
于是
\begin{equation}
  \xi(e_{ij})=q^j\frac{\partial}{\partial q^i}-p_i\frac{\partial}{\partial p_j}
  \tag{A.23}
  \label{eq:A.23}
\end{equation}
并且
\begin{equation}
  \mu(q^i,p_i;e_{ij})=-q^ip_j
  \tag{A.24}
  \label{eq:A.24}
\end{equation}

若进一步特殊化到 $SO(n)\subset GL(n,\mathbb{R})$、取基 $T_{ij}=e_{ij}-e_{ji}$，我们得到 $\mu(T_{ij})=q^jp_i-q^ip_j$；最后特殊化到 $T^{*}\mathbb{R}^3$、取 $T_i=\frac{1}{2}\epsilon_{ijk}T_{jk}$，我们发现 $\mu(T_i)=L_i$ 正是角动量的标准公式：
\begin{equation}
  \mu(q^i,p_i;T_i)=\epsilon_{ijk}q^jp_k
  \tag{A.25}
  \label{eq:A.25}
\end{equation}

\subsubsection{辛商} % 原文: A.1.2 Symplectic Quotient

若 $H^{1}_{DR}(M)=0$，可以证明
\begin{equation}
  \{\mu^A,\mu^B\}=f_C^{AB}\mu^C+c^{AB}
  \tag{A.26}
  \label{eq:A.26}
\end{equation}
其中 $c^{AB}$ 为某些常数。Poisson 括号上的 Jacobi 关系表明，这些 $c^{AB}$ 事实上定义了一个 Lie 代数 2-上循环（2-cocycle）。把 $\mu^A$ 平移常数，会把该上循环重新定义到相差一个上边缘（coboundary）。如果该 2-上循环为零，那么我们就能像在约束 Hamilton 动力学中那样同时施加约束 $\mu^A=0$。如果该 2-上循环代表一个非平凡的上同调类，那么「试图施加这些约束」就会遇到反常（anomaly）。

我们假定可以通过上边缘的平移把任何上循环 $c^{AB}$ 归零。若 $\zeta\in Z(\mathfrak{g})^{*}$ 落在中心（……的对偶）之中，则约束 $\phi=\mu-\zeta$ 是第一类约束（first class constraints）。也就是说，约束集合
\begin{equation}
  \phi^A=\mu(T^A)-\langle T^A,\zeta\rangle=0
  \tag{A.27}
  \label{eq:A.27}
\end{equation}
构成一组第一类约束。
% [Tutor Note: 原文此句为 "If $\zeta\in Z(\mathfrak g)^*$ is in the (dual of ) the center"，括号里明显漏了一个词或符号，原 PDF 上就是空白；译文照原样保留空缺。按上下文应读作「$\zeta$ 属于 $\mathfrak g$ 的中心 $Z(\mathfrak g)$ 的对偶」，即 $\zeta$ 只在中心上与生成元配对。这正是 $\phi^A$ 两两第一类（$\{\phi^A,\phi^B\}=f_C^{AB}\phi^C$）的条件，请对照 (A.26) 与 (A.28) 核实]

现在，取 $c^{AB}=0$ 的式 \eqref{eq:A.26} 给出
\begin{equation}
  \mathcal{L}_{V^A}(\mu^B)=f_C^{AB}\mu^C
  \tag{A.28}
  \label{eq:A.28}
\end{equation}
% [Tutor Note: (A.28) 左端的 $V^A$ 是 (A.4) 意义下 $\mu^A$ 的 Hamilton 向量场；这一步只是把 (A.26) 改写成 Lie 导数形式，$\mathcal{L}_{V^A}\mu^B$ 与 $\{\mu^A,\mu^B\}$ 至多差一个符号（取决于 $\iota(V)\omega=dh$ 的约定，见 (A.4) 处的批注）。若 $c^{AB}\neq0$，右端还要多出 $c^{AB}$，那正是「第二类约束／反常」的来源。另注意原文把结构常数排成 $f_C^{AB}$（上上／下），并非通常的 $f^{AB}{}_C$；引用时按 (A.26) 里与 $\mu^C$ 的缩并来配对，不要按上下标位置另设约定]

因此，若 $G$ 连通，我们就能断定：由 \eqref{eq:A.27} 所定义的相空间中的约束空间
\begin{equation}
  \mathcal{S}:=\{p\,|\,\phi^A(p)=0\}\subset M
  \tag{A.29}
  \label{eq:A.29}
\end{equation}
是 $G$-不变的。于是我们可以对 $G$-作用取商。

\textbf{定义} $M$ 关于 $G$ 的辛商（symplectic quotient）——或称 Hamilton 约化（Hamiltonian reduction）——是如下流形：
\begin{equation}
  M/\!/_{\zeta}G=\mu^{-1}(\zeta)/G
  \tag{A.30}
  \label{eq:A.30}
\end{equation}

\subsubsection{Kähler 商} % 原文: A.1.3 Kähler quotient

回忆：Kähler 流形是带有相容复结构 $I$ 的 Riemann 流形，即 $g(IX,IY)=g(X,Y)$（$X,Y\in TM$），且该复结构关于 Levi-Civita 联络平行，$\nabla I=0$。等价地，由 $\omega(X,Y)=g(X,IY)$ 定义的二形式是闭 2-形式。

用坐标来说：存在局部复坐标 $z^i$，使得 $ds^2=g_{i\bar{j}}dz^i\otimes d\bar{z}^{\bar{j}}+cc$，并且
\begin{equation}
  \omega=\frac{\mathrm{i}}{2}g_{i\bar{j}}dz^i\wedge d\bar{z}^{\bar{j}}
  \tag{A.31}
  \label{eq:A.31}
\end{equation}
是闭 2-形式。

当然，Kähler 流形就是带有辛形式 $\omega$ 的辛流形。

\textbf{定理}：设 $M$ 是 Kähler 流形，且 $G$ 无不动点地辛作用。那么辛商是光滑的 Kähler 流形，其 Kähler 结构继承自 $M$。

证明见所引文献。
% [Tutor Note: 作者以 "For a proof see the references" 跳过证明。这条定理是 Kempf–Ness 定理的 Kähler 版（即 Marsden–Weinstein 约化的复几何表述）：完整陈述还需要 $G$ 紧致、作用自由且 proper，原文的 "without fixed points" 只交代了自由性的一半；结论「辛商光滑并继承 Kähler 结构」依赖 $\mu^{-1}(\zeta)/G$ 与复商 $M^{\mathrm{stable}}/G$ 的同构，建议查 Kempf–Ness (1978) 或 Mumford《Geometric Invariant Theory》]

其一个重要特例如下。取 $V=\mathbb{C}^N$ 为复向量空间，辛形式为
\begin{equation}
  \omega=\frac{\mathrm{i}}{2}\sum_{a=1}^{N}dz^a\wedge d\bar{z}_a=\sum_a dx^a\wedge dy^a
  \tag{A.32}
  \label{eq:A.32}
\end{equation}
其中 $z^a=x^a+\sqrt{-1}\,y^a$。显然，这一形式在 $U(N)$ 作用于 $z^a$ 时保持不变。Lie 代数的一个元素是反 Hermite 矩阵
\begin{equation}
  (T^a{}_b)^{*}=-T^b{}_a
  \tag{A.33}
  \label{eq:A.33}
\end{equation}
相应的向量场是
\begin{equation}
  \xi(T)=z^bT^a{}_b\frac{\partial}{\partial z^a}-\bar{z}_b(T)^b{}_a\frac{\partial}{\partial\bar{z}_a}
  \tag{A.34}
  \label{eq:A.34}
\end{equation}
而相应的动量是
\begin{equation}
  \mu(T)=\mathrm{i}z^bT^a{}_b\bar{z}_a
  \tag{A.35}
  \label{eq:A.35}
\end{equation}
（相差一个常数）。

\textbf{例 1} $V=\mathbb{C}^N$，带有辛形式 \eqref{eq:A.32}，$G=U(1)$ 通过
\begin{equation}
  e^{\mathrm{i}\theta}:z^a\to e^{\mathrm{i}Q^a\theta}z^a,\qquad a=1,\dots,N
  \tag{A.36}
  \label{eq:A.36}
\end{equation}
作用。矩映射为
\begin{equation}
  \mu=\sum Q^a|z^a|^2
  \tag{A.37}
  \label{eq:A.37}
\end{equation}

辛商的性质强烈依赖于 $Q^a$ 以及矩映射的能级 $\zeta$：

\begin{enumerate}[label={\alph*.)},leftmargin=2.4em,itemsep=2pt]
  \item 除非 $Q^a$ 彼此可公度，$U(1)$ 作用是遍历的（ergodic）。让我们设 $Q^a$ 可公度。于是我们可以写 $Q^a=p^a/L$，其中 $L$ 为某个整数，且 $\mathrm{g.c.d}(p^1,\dots,p^N)=1$。商的性质依赖于整数 $p^i$。若它们两两互素，则商空间是光滑的。然而，假设存在子集 $S\subset\{1,\dots,N\}$，在其上 $p^a=rq^a$，其中 $r$ 为某个公因子。若对 $b\notin S$ 有 $z^b=0$，并取 $\theta=2\pi/r$，则该群有不动点。于是我们得到 $\mathbb{Z}_r$ 轨形奇点。
  \item 若一切 $Q^a$ 都为正且为有理数，并且 $\zeta>0$，则约束空间 $\mu=\zeta$ 是紧的且非空。它是一个形变球面。
  \item 若 $Q^a$ 符号不同，则约束空间非紧，并且可能出现非常不同的现象。
\end{enumerate}

这个例子也说明了 Kähler 商的拓扑如何随能级 $\zeta$ 的改变而改变。设有 $n_+$ 个 $Q_a=+1$ 与 $n_-$ 个 $Q_a=-1$。那么当 $\zeta>0$ 时辛商是丛 $\mathcal{O}(1)^{\oplus n_-}\to\mathbb{CP}^{n_+-1}$ 的全空间，而当 $\zeta<0$ 时它是 $\mathcal{O}(1)^{\oplus n_+}\to\mathbb{CP}^{n_--1}$。

\textbf{例 2} 超对称场论中的 $D$ 项。设我们有一个四维、带有四个超对称的超对称场论，其中有一组手征超场（chiral superfields）$\Phi^a$，它们按规范群的某个酉表示变换，该群的生成元为 $T^A$。设手征超场取值于一个线性空间，并具有标准的 Kähler 势 $\sum|\Phi^a|^2$。则势能中的 $D$ 项 $D^A$ 恰好是靶空间上 $G$-作用的矩映射。矩映射的能级 $\zeta^A$ 与 $G$ 中的 $U(1)$ 因子相联系，称为 Fayet–Iliopoulos 项。

{\small\itshape ［原文标记：*********************************］}
% [原文待办标记，原文 p.170 页末；原稿在此用一整行星号表示「此处还缺内容」，星号之后本片段结束]

\section{Kähler 量子化速览} % 原文: B. Lightning Review Of Kähler Quantization

\begin{keytakeaways}
  \item 出发点：标准相空间 $\mathcal{P}=\mathbb{R}^2$（坐标 $(q,p)$、辛形式 $\omega=dp\wedge dq/2\pi\hbar$，即 (B.1)），由此导致标准的 $[\hat{p},\hat{q}]=-\mathrm{i}\hbar$。
  \item 与通常量子化的差别：通常取 $\mathcal{H}=L^2(\mathbb{R})$，且为一般函数 $f(p,q)$ 定义 $Q(f)$ 必须额外选一套算符排序规定，以满足判据 (2.75)–(2.78)；Bargmann 表示改用 $\mathcal{H}=\mathrm{Hol}(\mathbb{C})$，即复平面上的整函数空间。
  \item 核心构造：内积带 Gaussian 权 $e^{-\kappa|z|^2}$ (B.2)，$\kappa>0$ 扮演 $1/\hbar$；一般（不必全纯的）可积函数 $f(z,\bar{z})$ 的量子化由积分核 (B.4) 给出，其算符形式为 $Q^{(\kappa)}(f)=\Pi\circ M_f\circ\iota$。
  \item 代数实现：取 $a=\sqrt{\kappa}\,Q^{(\kappa)}(\bar{z})$、$a^\dagger=\sqrt{\kappa}\,Q^{(\kappa)}(z)$ (B.5)，则 $[a,a^\dagger]=1$ (B.6)——即 $a$ 与 $a^\dagger$ 分别由「乘以 $\bar{z}$ 后投影」与「乘以 $z$ 后投影」实现。
\end{keytakeaways}

\subsection{Bargmann 表示} % 原文: B.1 The Bargmann Representation

Bargmann 表示是量子化标准相空间 $\mathcal{P}=\mathbb{R}^2$ 的另一种方式；该相空间的坐标为 $(q,p)$，辛形式为
\begin{equation}
  \omega=\frac{dp\wedge dq}{2\pi\hbar}
  \tag{B.1}
  \label{eq:B.1}
\end{equation}
由此得到标准的 $[\hat{p},\hat{q}]=-\mathrm{i}\hbar$。{\small\itshape ［原文标记：$\clubsuit$符号？？？$\clubsuit$］}
% [原文待办标记，原文 p.171 页边；作者自己怀疑 (B.1) 与对易子的符号约定]

在通常的量子化中，我们把 $\mathcal{P}=T^{*}\mathbb{R}$ 看作带有实轴上坐标 $q$ 的相空间，于是 $\mathcal{H}=L^2(\mathbb{R})$，波函数是可归一化的函数 $\psi(q)$，而 $\hat{p}=-\frac{\mathrm{i}}{\hbar}\frac{d}{dq}$，等等。为了对一般函数 $f(p,q)$ 定义 $Q^{(k)}(f)$，必须选取一套算符排序（operator ordering）规定；这可以用若干种方式做到，以满足上文所列的判据 \eqref{eq:2.75}–\eqref{eq:2.78}。

在 Bargmann 表示中，Hilbert 空间则取为 $\mathcal{H}=\mathrm{Hol}(\mathbb{C})$，即复平面上全体整函数构成的空间，其在如下定义的内积下可归一化：
\begin{equation}
  \langle\psi_1,\psi_2\rangle=\int_{\mathbb{C}}e^{-\kappa|z|^2}\psi_1(z)^{*}\psi_2(z)\,\omega
  \tag{B.2}
  \label{eq:B.2}
\end{equation}

这里 $\kappa>0$ 将扮演 $1/\hbar$ 的角色，而
\begin{equation}
  \omega=\frac{\mathrm{i}\,dz\wedge d\bar{z}}{2\pi}=\frac{dx\,dy}{\pi}=\frac{d\phi}{2\pi}d(r^2)
  \tag{B.3}
  \label{eq:B.3}
\end{equation}

可积函数 $f(z,\bar{z})$（\underline{不}必全纯或反全纯）的量子化 $Q^{(k)}(f)$ 由下式定义
\begin{equation}
  \langle\psi_{j_1}|Q^{(\kappa)}(f)|\psi_{j_2}\rangle=\int_{\mathbb{C}}e^{-\kappa|z|^2}\psi_{j_1}^{*}(\bar{z})f(z,\bar{z})\psi_{j_2}(z)\,\frac{dx\,dy}{\pi}
  \tag{B.4}
  \label{eq:B.4}
\end{equation}
% [Tutor Note: (B.4) 印刷为 $\psi^{*}_{j_1}(\bar{z})$，但 $\psi_{j_1}\in\mathrm{Hol}(\mathbb{C})$ 是全纯函数，按 (B.2) 的内积此处应是 $\psi^{*}_{j_1}(z)$；页面图像上确实印着 $\bar z$，故照录。同一页还有 $Q^{(k)}$ 与 $Q^{(\kappa)}$ 混用（(B.4) 前的正文写 $Q^{(k)}$，而 (B.4)、(B.5) 写 $Q^{(\kappa)}$），二者指同一个参数，即 (B.2) 中扮演 $1/\hbar$ 的 $\kappa$]

特别地，若我们定义
\begin{equation}
  a=\sqrt{\kappa}\,Q^{(\kappa)}(\bar{z}),\qquad a^{\dagger}=\sqrt{\kappa}\,Q^{(\kappa)}(z)
  \tag{B.5}
  \label{eq:B.5}
\end{equation}
那么我们就能验证通常的关系
\begin{equation}
  [a,a^{\dagger}]=1
  \tag{B.6}
  \label{eq:B.6}
\end{equation}

{\small\itshape ［原文标记：*********************************\newline
关于相干态的一些评注。它们为何对半经典计算有用。\newline
BEREZIN 变换。\newline
*********************************］}
% [原文待办标记，原文 p.171 页末；原稿在两行星号之间列出了三点待补内容（相干态、半经典计算的用处、Berezin 变换），此处照原样保留]
